% =====================================================================
%  Factor construction formulas.
%
%  GENERATED by scripts/render_formula_appendix.py from
%  backend/pipeline/factor_defs.py via backend/pipeline/formula.py.
%  Do not edit: rerun the script after changing a factor definition.
%
%  Included by factor-model-paper.tex.
% =====================================================================
\newpage
\section{Factor construction formulas}
\label{app:formulas}

Every factor in the model, written out: the arithmetic that builds it from named
instruments, and the equation that residualises it against the blocks above it.
Section \ref{sec:data} gives the reasoning behind the constructions; this states
them.

This appendix is generated from the factor registry rather than transcribed, so
it cannot drift from the code that builds the panel. The transcription from the
builders is itself checked: each rendered formula is evaluated independently by
the test suite and asserted to reproduce the builder's own output on simulated
inputs, which is what makes an equation here a statement about the
implementation rather than about an intention.

\subsection{Reading these}

Each factor is stored twice, and the two are different series. The raw column
\code{fact\_factor\_return.ret\_excess} holds $f_t$, the construction formula's
own output. The orthogonalised column \code{fact\_factor\_return.ret\_orth}
holds $\tilde{f}_t$, after the block hierarchy has residualised it. The model
may be estimated on either, and the choice is recorded on the spec and carried
through to the covariance, because $\beta'\Sigma\beta$ requires $\Sigma$ to be
the covariance of the same series the $\beta$s refer to. A factor whose
orthogonalisation list is empty has $\tilde{f}_t = f_t$ identically.

\begin{center}
\small
\begin{tabular}{@{}ll@{}}
\toprule
Symbol & Meaning \\
\midrule
$r^X_t$ & log total return of instrument $X$, on a price series already adjusted
          for dividends and splits \\
$c_t$ & daily cash rate: \code{FRED:DFF} at $t$, annualised percent, divided by
        100 and by 252 \\
$x^X_t$ & excess return over cash, $x^X_t = r^X_t - c_{t-1}$ \\
$\Delta$ & first difference, $\Delta s_t = s_t - s_{t-1}$ \\
$f_t$ & the raw factor \\
$\tilde{f}_t$ & the orthogonalised factor \\
$\tau$ & the most recent orthogonalisation refit at or before $t$ \\
\bottomrule
\end{tabular}
\end{center}

Every lag is deliberate. A quantity dated $t-1$ is there because using its value
at $t$ would put information into a return that was not available when the
return was earned.


% ---------------------------------------------------
\subsection{Equity Market}

\subsubsection{\code{eq\_global} --- Global Equity}
Hierarchy level 0. Method \code{single}.

\begin{quote}\small MSCI ACWI total return less cash. The single global market factor.\end{quote}

\noindent\textbf{Construction}

\eqn{f_t = x^{\text{ACWI}}_t}

\begin{enumerate}[leftmargin=*,itemsep=2pt]
  \item Subtract the overnight rate: x\_t = r\_t - c\_(t-1). ACWI is a funded long position, so its return has to be measured over cash.
\end{enumerate}

\begin{center}\small
\begin{tabularx}{\linewidth}{@{}llX@{}}
\toprule
Symbol & Source & Meaning \\
\midrule
$r^{\text{ACWI}}_t$ & \code{ACWI} & instrument return: log total return of ACWI \\
$c_t$ & \code{FRED:DFF} & daily cash rate: FRED:DFF\_t / 100 / 252, lagged one day because the overnight rate earned on day t is set at t-1 \\
\bottomrule
\end{tabularx}
\end{center}

\noindent\textbf{Orthogonalisation}

\eqn{\tilde{f}_t = f_t}

Nothing sits above this factor in the hierarchy, so the raw and orthogonalised panels hold the same series.

\subsubsection{\code{eq\_us} --- US Equity (ex-global)}
Hierarchy level 1. Method \code{single}.

\noindent\textbf{Construction}

\eqn{f_t = x^{\text{IVV}}_t}

\begin{enumerate}[leftmargin=*,itemsep=2pt]
  \item Subtract the overnight rate: x\_t = r\_t - c\_(t-1). IVV is a funded long position, so its return has to be measured over cash.
\end{enumerate}

\begin{center}\small
\begin{tabularx}{\linewidth}{@{}llX@{}}
\toprule
Symbol & Source & Meaning \\
\midrule
$r^{\text{IVV}}_t$ & \code{IVV} & instrument return: log total return of IVV \\
$c_t$ & \code{FRED:DFF} & daily cash rate: FRED:DFF\_t / 100 / 252, lagged one day because the overnight rate earned on day t is set at t-1 \\
\bottomrule
\end{tabularx}
\end{center}

\noindent\textbf{Orthogonalisation}

\eqn{\tilde{f}_t = f_t - \hat{a}_\tau - \hat{b}_{1,\tau}\,\tilde{f}^{\text{eq\_global}}_t}

\eqn{(\hat{a}_\tau, \hat{b}_\tau) = \arg\min_{a,b} \sum_{s=\tau-504}^{\tau-1} \Bigl(f_s - a - \sum_j b_j\,\tilde{f}^{g_j}_s\Bigr)^2}

with $g_1$ = \code{eq\_global}, each taken as its own \emph{orthogonalised} series, and $\tau$ the most recent refit at or before $t$.

\subsubsection{\code{eq\_europe} --- Europe Equity (ex-global)}
Hierarchy level 1. Method \code{single}.

\noindent\textbf{Construction}

\eqn{f_t = x^{\text{VGK}}_t}

\begin{enumerate}[leftmargin=*,itemsep=2pt]
  \item Subtract the overnight rate: x\_t = r\_t - c\_(t-1). VGK is a funded long position, so its return has to be measured over cash.
\end{enumerate}

\begin{center}\small
\begin{tabularx}{\linewidth}{@{}llX@{}}
\toprule
Symbol & Source & Meaning \\
\midrule
$r^{\text{VGK}}_t$ & \code{VGK} & instrument return: log total return of VGK \\
$c_t$ & \code{FRED:DFF} & daily cash rate: FRED:DFF\_t / 100 / 252, lagged one day because the overnight rate earned on day t is set at t-1 \\
\bottomrule
\end{tabularx}
\end{center}

\noindent\textbf{Orthogonalisation}

\eqn{\tilde{f}_t = f_t - \hat{a}_\tau - \hat{b}_{1,\tau}\,\tilde{f}^{\text{eq\_global}}_t}

\eqn{(\hat{a}_\tau, \hat{b}_\tau) = \arg\min_{a,b} \sum_{s=\tau-504}^{\tau-1} \Bigl(f_s - a - \sum_j b_j\,\tilde{f}^{g_j}_s\Bigr)^2}

with $g_1$ = \code{eq\_global}, each taken as its own \emph{orthogonalised} series, and $\tau$ the most recent refit at or before $t$.

\subsubsection{\code{eq\_japan} --- Japan Equity (ex-global)}
Hierarchy level 1. Method \code{single}.

\noindent\textbf{Construction}

\eqn{f_t = x^{\text{EWJ}}_t}

\begin{enumerate}[leftmargin=*,itemsep=2pt]
  \item Subtract the overnight rate: x\_t = r\_t - c\_(t-1). EWJ is a funded long position, so its return has to be measured over cash.
\end{enumerate}

\begin{center}\small
\begin{tabularx}{\linewidth}{@{}llX@{}}
\toprule
Symbol & Source & Meaning \\
\midrule
$r^{\text{EWJ}}_t$ & \code{EWJ} & instrument return: log total return of EWJ \\
$c_t$ & \code{FRED:DFF} & daily cash rate: FRED:DFF\_t / 100 / 252, lagged one day because the overnight rate earned on day t is set at t-1 \\
\bottomrule
\end{tabularx}
\end{center}

\noindent\textbf{Orthogonalisation}

\eqn{\tilde{f}_t = f_t - \hat{a}_\tau - \hat{b}_{1,\tau}\,\tilde{f}^{\text{eq\_global}}_t}

\eqn{(\hat{a}_\tau, \hat{b}_\tau) = \arg\min_{a,b} \sum_{s=\tau-504}^{\tau-1} \Bigl(f_s - a - \sum_j b_j\,\tilde{f}^{g_j}_s\Bigr)^2}

with $g_1$ = \code{eq\_global}, each taken as its own \emph{orthogonalised} series, and $\tau$ the most recent refit at or before $t$.

\subsubsection{\code{eq\_em} --- EM Equity (ex-global)}
Hierarchy level 1. Method \code{single}.

\noindent\textbf{Construction}

\eqn{f_t = x^{\text{EEM}}_t}

\begin{enumerate}[leftmargin=*,itemsep=2pt]
  \item Subtract the overnight rate: x\_t = r\_t - c\_(t-1). EEM is a funded long position, so its return has to be measured over cash.
\end{enumerate}

\begin{center}\small
\begin{tabularx}{\linewidth}{@{}llX@{}}
\toprule
Symbol & Source & Meaning \\
\midrule
$r^{\text{EEM}}_t$ & \code{EEM} & instrument return: log total return of EEM \\
$c_t$ & \code{FRED:DFF} & daily cash rate: FRED:DFF\_t / 100 / 252, lagged one day because the overnight rate earned on day t is set at t-1 \\
\bottomrule
\end{tabularx}
\end{center}

\noindent\textbf{Orthogonalisation}

\eqn{\tilde{f}_t = f_t - \hat{a}_\tau - \hat{b}_{1,\tau}\,\tilde{f}^{\text{eq\_global}}_t}

\eqn{(\hat{a}_\tau, \hat{b}_\tau) = \arg\min_{a,b} \sum_{s=\tau-504}^{\tau-1} \Bigl(f_s - a - \sum_j b_j\,\tilde{f}^{g_j}_s\Bigr)^2}

with $g_1$ = \code{eq\_global}, each taken as its own \emph{orthogonalised} series, and $\tau$ the most recent refit at or before $t$.

\subsubsection{\code{eq\_size} --- Size (small minus large)}
Hierarchy level 1. Method \code{spread}.

\begin{quote}\small Russell 2000 minus Russell 1000. Self-financing, so no cash leg.\end{quote}

\noindent\textbf{Construction}

\eqn{f_t = r^{\text{IWM}}_t - r^{\text{IWB}}_t}

\begin{enumerate}[leftmargin=*,itemsep=2pt]
  \item Self-financing, so no cash leg: the funding cost cancels between the two legs and subtracting it would understate the spread.
\end{enumerate}

\begin{center}\small
\begin{tabularx}{\linewidth}{@{}llX@{}}
\toprule
Symbol & Source & Meaning \\
\midrule
$r^{\text{IWM}}_t$ & \code{IWM} & long leg: log total return of IWM \\
$r^{\text{IWB}}_t$ & \code{IWB} & short leg: log total return of IWB \\
\bottomrule
\end{tabularx}
\end{center}

\noindent\textbf{Orthogonalisation}

\eqn{\tilde{f}_t = f_t - \hat{a}_\tau - \hat{b}_{1,\tau}\,\tilde{f}^{\text{eq\_global}}_t}

\eqn{(\hat{a}_\tau, \hat{b}_\tau) = \arg\min_{a,b} \sum_{s=\tau-504}^{\tau-1} \Bigl(f_s - a - \sum_j b_j\,\tilde{f}^{g_j}_s\Bigr)^2}

with $g_1$ = \code{eq\_global}, each taken as its own \emph{orthogonalised} series, and $\tau$ the most recent refit at or before $t$.

\subsubsection{\code{eq\_cyclical} --- Cyclical minus Defensive}
Hierarchy level 1. Method \code{basket}.

\begin{quote}\small Sector breadth without spending eleven factors on it.\end{quote}

\noindent\textbf{Construction}

\eqn{f_t = \frac{r^{\text{XLI}}_t + r^{\text{XLB}}_t + r^{\text{XLE}}_t + r^{\text{XLF}}_t + r^{\text{XLY}}_t}{5} - \frac{r^{\text{XLP}}_t + r^{\text{XLU}}_t + r^{\text{XLV}}_t}{3}}

\begin{enumerate}[leftmargin=*,itemsep=2pt]
  \item Equally weighted long basket of 5: each leg carries 1/5 of the notional.
  \item Minus an equally weighted short basket of 3. Both sides carry the market direction, so it cancels and what survives is the spread between them.
  \item Self-financing, so no cash leg.
\end{enumerate}

\begin{center}\small
\begin{tabularx}{\linewidth}{@{}llX@{}}
\toprule
Symbol & Source & Meaning \\
\midrule
$r^{\text{XLI}}_t$ & \code{XLI} & long leg: log total return of XLI \\
$r^{\text{XLB}}_t$ & \code{XLB} & long leg: log total return of XLB \\
$r^{\text{XLE}}_t$ & \code{XLE} & long leg: log total return of XLE \\
$r^{\text{XLF}}_t$ & \code{XLF} & long leg: log total return of XLF \\
$r^{\text{XLY}}_t$ & \code{XLY} & long leg: log total return of XLY \\
$r^{\text{XLP}}_t$ & \code{XLP} & short leg: log total return of XLP \\
$r^{\text{XLU}}_t$ & \code{XLU} & short leg: log total return of XLU \\
$r^{\text{XLV}}_t$ & \code{XLV} & short leg: log total return of XLV \\
\bottomrule
\end{tabularx}
\end{center}

\noindent\textbf{Orthogonalisation}

\eqn{\tilde{f}_t = f_t - \hat{a}_\tau - \hat{b}_{1,\tau}\,\tilde{f}^{\text{eq\_global}}_t}

\eqn{(\hat{a}_\tau, \hat{b}_\tau) = \arg\min_{a,b} \sum_{s=\tau-504}^{\tau-1} \Bigl(f_s - a - \sum_j b_j\,\tilde{f}^{g_j}_s\Bigr)^2}

with $g_1$ = \code{eq\_global}, each taken as its own \emph{orthogonalised} series, and $\tau$ the most recent refit at or before $t$.


% ---------------------------------------------------
\subsection{Equity Style}

\subsubsection{\code{sty\_value} --- Value}
Hierarchy level 2. Method \code{single}.

\begin{quote}\small Long-only ETF residualised to market and sector. An approximation to a true long-short HML; validated against F-F HML over the common sample.\end{quote}

\noindent\textbf{Construction}

\eqn{f_t = x^{\text{VLUE}}_t}

\begin{enumerate}[leftmargin=*,itemsep=2pt]
  \item Subtract the overnight rate: x\_t = r\_t - c\_(t-1). VLUE is a funded long position, so its return has to be measured over cash.
\end{enumerate}

\begin{center}\small
\begin{tabularx}{\linewidth}{@{}llX@{}}
\toprule
Symbol & Source & Meaning \\
\midrule
$r^{\text{VLUE}}_t$ & \code{VLUE} & instrument return: log total return of VLUE \\
$c_t$ & \code{FRED:DFF} & daily cash rate: FRED:DFF\_t / 100 / 252, lagged one day because the overnight rate earned on day t is set at t-1 \\
\bottomrule
\end{tabularx}
\end{center}

\noindent\textbf{Orthogonalisation}

\eqn{\tilde{f}_t = f_t - \hat{a}_\tau - \hat{b}_{1,\tau}\,\tilde{f}^{\text{eq\_global}}_t - \hat{b}_{2,\tau}\,\tilde{f}^{\text{eq\_us}}_t - \hat{b}_{3,\tau}\,\tilde{f}^{\text{eq\_cyclical}}_t}

\eqn{(\hat{a}_\tau, \hat{b}_\tau) = \arg\min_{a,b} \sum_{s=\tau-504}^{\tau-1} \Bigl(f_s - a - \sum_j b_j\,\tilde{f}^{g_j}_s\Bigr)^2}

with $g_1, \dots, g_{3}$ = \code{eq\_global}, \code{eq\_us}, \code{eq\_cyclical}, each taken as its own \emph{orthogonalised} series, and $\tau$ the most recent refit at or before $t$.

\subsubsection{\code{sty\_momentum} --- Momentum}
Hierarchy level 2. Method \code{single}.

\noindent\textbf{Construction}

\eqn{f_t = x^{\text{MTUM}}_t}

\begin{enumerate}[leftmargin=*,itemsep=2pt]
  \item Subtract the overnight rate: x\_t = r\_t - c\_(t-1). MTUM is a funded long position, so its return has to be measured over cash.
\end{enumerate}

\begin{center}\small
\begin{tabularx}{\linewidth}{@{}llX@{}}
\toprule
Symbol & Source & Meaning \\
\midrule
$r^{\text{MTUM}}_t$ & \code{MTUM} & instrument return: log total return of MTUM \\
$c_t$ & \code{FRED:DFF} & daily cash rate: FRED:DFF\_t / 100 / 252, lagged one day because the overnight rate earned on day t is set at t-1 \\
\bottomrule
\end{tabularx}
\end{center}

\noindent\textbf{Orthogonalisation}

\eqn{\tilde{f}_t = f_t - \hat{a}_\tau - \hat{b}_{1,\tau}\,\tilde{f}^{\text{eq\_global}}_t - \hat{b}_{2,\tau}\,\tilde{f}^{\text{eq\_us}}_t - \hat{b}_{3,\tau}\,\tilde{f}^{\text{eq\_cyclical}}_t}

\eqn{(\hat{a}_\tau, \hat{b}_\tau) = \arg\min_{a,b} \sum_{s=\tau-504}^{\tau-1} \Bigl(f_s - a - \sum_j b_j\,\tilde{f}^{g_j}_s\Bigr)^2}

with $g_1, \dots, g_{3}$ = \code{eq\_global}, \code{eq\_us}, \code{eq\_cyclical}, each taken as its own \emph{orthogonalised} series, and $\tau$ the most recent refit at or before $t$.

\subsubsection{\code{sty\_quality} --- Quality}
Hierarchy level 2. Method \code{single}.

\noindent\textbf{Construction}

\eqn{f_t = x^{\text{QUAL}}_t}

\begin{enumerate}[leftmargin=*,itemsep=2pt]
  \item Subtract the overnight rate: x\_t = r\_t - c\_(t-1). QUAL is a funded long position, so its return has to be measured over cash.
\end{enumerate}

\begin{center}\small
\begin{tabularx}{\linewidth}{@{}llX@{}}
\toprule
Symbol & Source & Meaning \\
\midrule
$r^{\text{QUAL}}_t$ & \code{QUAL} & instrument return: log total return of QUAL \\
$c_t$ & \code{FRED:DFF} & daily cash rate: FRED:DFF\_t / 100 / 252, lagged one day because the overnight rate earned on day t is set at t-1 \\
\bottomrule
\end{tabularx}
\end{center}

\noindent\textbf{Orthogonalisation}

\eqn{\tilde{f}_t = f_t - \hat{a}_\tau - \hat{b}_{1,\tau}\,\tilde{f}^{\text{eq\_global}}_t - \hat{b}_{2,\tau}\,\tilde{f}^{\text{eq\_us}}_t - \hat{b}_{3,\tau}\,\tilde{f}^{\text{eq\_cyclical}}_t}

\eqn{(\hat{a}_\tau, \hat{b}_\tau) = \arg\min_{a,b} \sum_{s=\tau-504}^{\tau-1} \Bigl(f_s - a - \sum_j b_j\,\tilde{f}^{g_j}_s\Bigr)^2}

with $g_1, \dots, g_{3}$ = \code{eq\_global}, \code{eq\_us}, \code{eq\_cyclical}, each taken as its own \emph{orthogonalised} series, and $\tau$ the most recent refit at or before $t$.

\subsubsection{\code{sty\_lowvol} --- Low Volatility}
Hierarchy level 2. Method \code{single}.

\noindent\textbf{Construction}

\eqn{f_t = x^{\text{USMV}}_t}

\begin{enumerate}[leftmargin=*,itemsep=2pt]
  \item Subtract the overnight rate: x\_t = r\_t - c\_(t-1). USMV is a funded long position, so its return has to be measured over cash.
\end{enumerate}

\begin{center}\small
\begin{tabularx}{\linewidth}{@{}llX@{}}
\toprule
Symbol & Source & Meaning \\
\midrule
$r^{\text{USMV}}_t$ & \code{USMV} & instrument return: log total return of USMV \\
$c_t$ & \code{FRED:DFF} & daily cash rate: FRED:DFF\_t / 100 / 252, lagged one day because the overnight rate earned on day t is set at t-1 \\
\bottomrule
\end{tabularx}
\end{center}

\noindent\textbf{Orthogonalisation}

\eqn{\tilde{f}_t = f_t - \hat{a}_\tau - \hat{b}_{1,\tau}\,\tilde{f}^{\text{eq\_global}}_t - \hat{b}_{2,\tau}\,\tilde{f}^{\text{eq\_us}}_t - \hat{b}_{3,\tau}\,\tilde{f}^{\text{eq\_cyclical}}_t}

\eqn{(\hat{a}_\tau, \hat{b}_\tau) = \arg\min_{a,b} \sum_{s=\tau-504}^{\tau-1} \Bigl(f_s - a - \sum_j b_j\,\tilde{f}^{g_j}_s\Bigr)^2}

with $g_1, \dots, g_{3}$ = \code{eq\_global}, \code{eq\_us}, \code{eq\_cyclical}, each taken as its own \emph{orthogonalised} series, and $\tau$ the most recent refit at or before $t$.


% ---------------------------------------------------
\subsection{Rates}

\subsubsection{\code{rt\_us\_level} --- US Rates Level}
Hierarchy level 0. Method \code{curve}.

\begin{quote}\small Equal unit-duration long across the curve, scaled to 5y duration.\end{quote}

\noindent\textbf{Construction}

\eqn{f_t = 5 \cdot \frac{1}{4}\sum_{T} u^{(T)}_t, \qquad u^{(T)}_t = \frac{b^{(T)}_t}{D^{(T)}_{t-1}}}

\eqn{b^{(T)}_t = -D^{(T)}_{t-1}\,\Delta y^{(T)}_t + \tfrac{1}{2} C^{(T)}_{t-1}\,\bigl(\Delta y^{(T)}_t\bigr)^2 + \frac{y^{(T)}_{t-1} - c_{t-1}}{252}}

\begin{enumerate}[leftmargin=*,itemsep=2pt]
  \item Convert each tenor's yield change into a bond return, with duration and convexity evaluated at yesterday's yield. Section 2.2 requires every factor to be a return, and a yield change is not one.
  \item Divide each leg by its own modified duration. Five basis points on a 30y is not the same trade as five on a 2y, and without this the long end would dominate every shape.
  \item Level: an equal unit-duration long at every tenor (2y, 5y, 10y, 30y).
  \item Scale back by 5 years of duration.
\end{enumerate}

\begin{center}\small
\begin{tabularx}{\linewidth}{@{}llX@{}}
\toprule
Symbol & Source & Meaning \\
\midrule
$c_t$ & \code{FRED:DFF} & daily cash rate: FRED:DFF\_t / 100 / 252, lagged one day because the overnight rate earned on day t is set at t-1 \\
$\text{FRED:DGS2}$ & \code{FRED:DGS2} & US\_TSY par yield at 2y \\
$\text{FRED:DGS5}$ & \code{FRED:DGS5} & US\_TSY par yield at 5y \\
$\text{FRED:DGS10}$ & \code{FRED:DGS10} & US\_TSY par yield at 10y \\
$\text{FRED:DGS30}$ & \code{FRED:DGS30} & US\_TSY par yield at 30y \\
\bottomrule
\end{tabularx}
\end{center}

\noindent\textbf{Orthogonalisation}

\eqn{\tilde{f}_t = f_t}

Nothing sits above this factor in the hierarchy, so the raw and orthogonalised panels hold the same series.

\subsubsection{\code{rt\_us\_slope} --- US Rates Slope (10s-2s)}
Hierarchy level 0. Method \code{curve}.

\begin{quote}\small Duration-neutral steepener: long 10y, short 2y, equal duration.\end{quote}

\noindent\textbf{Construction}

\eqn{f_t = 5 \cdot \bigl(u^{(10y)}_t - u^{(2y)}_t\bigr), \qquad u^{(T)}_t = \frac{b^{(T)}_t}{D^{(T)}_{t-1}}}

\eqn{b^{(T)}_t = -D^{(T)}_{t-1}\,\Delta y^{(T)}_t + \tfrac{1}{2} C^{(T)}_{t-1}\,\bigl(\Delta y^{(T)}_t\bigr)^2 + \frac{y^{(T)}_{t-1} - c_{t-1}}{252}}

\begin{enumerate}[leftmargin=*,itemsep=2pt]
  \item Convert each tenor's yield change into a bond return, with duration and convexity evaluated at yesterday's yield. Section 2.2 requires every factor to be a return, and a yield change is not one.
  \item Divide each leg by its own modified duration. Five basis points on a 30y is not the same trade as five on a 2y, and without this the long end would dominate every shape.
  \item Slope: long 10y against short 2y in equal units of duration, so a parallel shift nets to zero and only the steepening survives.
  \item Scale back by 5 years of duration.
\end{enumerate}

\begin{center}\small
\begin{tabularx}{\linewidth}{@{}llX@{}}
\toprule
Symbol & Source & Meaning \\
\midrule
$c_t$ & \code{FRED:DFF} & daily cash rate: FRED:DFF\_t / 100 / 252, lagged one day because the overnight rate earned on day t is set at t-1 \\
$\text{FRED:DGS2}$ & \code{FRED:DGS2} & US\_TSY par yield at 2y \\
$\text{FRED:DGS10}$ & \code{FRED:DGS10} & US\_TSY par yield at 10y \\
\bottomrule
\end{tabularx}
\end{center}

\noindent\textbf{Orthogonalisation}

\eqn{\tilde{f}_t = f_t}

Nothing sits above this factor in the hierarchy, so the raw and orthogonalised panels hold the same series.

\subsubsection{\code{rt\_us\_curve} --- US Rates Curvature (butterfly)}
Hierarchy level 0. Method \code{curve}.

\begin{quote}\small Long the belly, short the wings, duration-neutral.\end{quote}

\noindent\textbf{Construction}

\eqn{f_t = 5 \cdot \bigl(2u^{(5y)}_t - u^{(2y)}_t - u^{(30y)}_t\bigr), \qquad u^{(T)}_t = \frac{b^{(T)}_t}{D^{(T)}_{t-1}}}

\eqn{b^{(T)}_t = -D^{(T)}_{t-1}\,\Delta y^{(T)}_t + \tfrac{1}{2} C^{(T)}_{t-1}\,\bigl(\Delta y^{(T)}_t\bigr)^2 + \frac{y^{(T)}_{t-1} - c_{t-1}}{252}}

\begin{enumerate}[leftmargin=*,itemsep=2pt]
  \item Convert each tenor's yield change into a bond return, with duration and convexity evaluated at yesterday's yield. Section 2.2 requires every factor to be a return, and a yield change is not one.
  \item Divide each leg by its own modified duration. Five basis points on a 30y is not the same trade as five on a 2y, and without this the long end would dominate every shape.
  \item Curvature: long two units of the 5y belly against one unit each of the 2y and 30y wings. Neutral to both a parallel shift and a steepening.
  \item Scale back by 5 years of duration.
\end{enumerate}

\begin{center}\small
\begin{tabularx}{\linewidth}{@{}llX@{}}
\toprule
Symbol & Source & Meaning \\
\midrule
$c_t$ & \code{FRED:DFF} & daily cash rate: FRED:DFF\_t / 100 / 252, lagged one day because the overnight rate earned on day t is set at t-1 \\
$\text{FRED:DGS2}$ & \code{FRED:DGS2} & US\_TSY par yield at 2y \\
$\text{FRED:DGS5}$ & \code{FRED:DGS5} & US\_TSY par yield at 5y \\
$\text{FRED:DGS30}$ & \code{FRED:DGS30} & US\_TSY par yield at 30y \\
\bottomrule
\end{tabularx}
\end{center}

\noindent\textbf{Orthogonalisation}

\eqn{\tilde{f}_t = f_t}

Nothing sits above this factor in the hierarchy, so the raw and orthogonalised panels hold the same series.

\subsubsection{\code{rt\_ea\_level} --- EA Rates Level}
Hierarchy level 0. Method \code{curve}.

\begin{quote}\small ECB AAA curve. Warehouse-sourced, so it lags the US block.\end{quote}

\noindent\textbf{Construction}

\eqn{f_t = 5 \cdot \frac{1}{4}\sum_{T} u^{(T)}_t, \qquad u^{(T)}_t = \frac{b^{(T)}_t}{D^{(T)}_{t-1}}}

\eqn{b^{(T)}_t = -D^{(T)}_{t-1}\,\Delta y^{(T)}_t + \tfrac{1}{2} C^{(T)}_{t-1}\,\bigl(\Delta y^{(T)}_t\bigr)^2 + \frac{y^{(T)}_{t-1} - c_{t-1}}{252}}

\begin{enumerate}[leftmargin=*,itemsep=2pt]
  \item Convert each tenor's yield change into a bond return, with duration and convexity evaluated at yesterday's yield. Section 2.2 requires every factor to be a return, and a yield change is not one.
  \item Divide each leg by its own modified duration. Five basis points on a 30y is not the same trade as five on a 2y, and without this the long end would dominate every shape.
  \item Level: an equal unit-duration long at every tenor (2y, 5y, 10y, 30y).
  \item Scale back by 5 years of duration.
\end{enumerate}

\begin{center}\small
\begin{tabularx}{\linewidth}{@{}llX@{}}
\toprule
Symbol & Source & Meaning \\
\midrule
$c_t$ & \code{FRED:DFF} & daily cash rate: FRED:DFF\_t / 100 / 252, lagged one day because the overnight rate earned on day t is set at t-1 \\
$\text{ECB:BUND\_2Y}$ & \code{ECB:BUND\_2Y} & EA\_AAA par yield at 2y \\
$\text{ECB:BUND\_5Y}$ & \code{ECB:BUND\_5Y} & EA\_AAA par yield at 5y \\
$\text{ECB:BUND\_10Y}$ & \code{ECB:BUND\_10Y} & EA\_AAA par yield at 10y \\
$\text{ECB:BUND\_30Y}$ & \code{ECB:BUND\_30Y} & EA\_AAA par yield at 30y \\
\bottomrule
\end{tabularx}
\end{center}

\noindent\textbf{Orthogonalisation}

\eqn{\tilde{f}_t = f_t - \hat{a}_\tau - \hat{b}_{1,\tau}\,\tilde{f}^{\text{rt\_us\_level}}_t}

\eqn{(\hat{a}_\tau, \hat{b}_\tau) = \arg\min_{a,b} \sum_{s=\tau-504}^{\tau-1} \Bigl(f_s - a - \sum_j b_j\,\tilde{f}^{g_j}_s\Bigr)^2}

with $g_1$ = \code{rt\_us\_level}, each taken as its own \emph{orthogonalised} series, and $\tau$ the most recent refit at or before $t$.

\subsubsection{\code{rt\_jp\_level} --- JP Rates Level}
Hierarchy level 0. Method \code{curve}.

\begin{quote}\small MOF/BOJ curve. Warehouse-sourced, so it lags the US block.\end{quote}

\noindent\textbf{Construction}

\eqn{f_t = 5 \cdot \frac{1}{4}\sum_{T} u^{(T)}_t, \qquad u^{(T)}_t = \frac{b^{(T)}_t}{D^{(T)}_{t-1}}}

\eqn{b^{(T)}_t = -D^{(T)}_{t-1}\,\Delta y^{(T)}_t + \tfrac{1}{2} C^{(T)}_{t-1}\,\bigl(\Delta y^{(T)}_t\bigr)^2 + \frac{y^{(T)}_{t-1} - c_{t-1}}{252}}

\begin{enumerate}[leftmargin=*,itemsep=2pt]
  \item Convert each tenor's yield change into a bond return, with duration and convexity evaluated at yesterday's yield. Section 2.2 requires every factor to be a return, and a yield change is not one.
  \item Divide each leg by its own modified duration. Five basis points on a 30y is not the same trade as five on a 2y, and without this the long end would dominate every shape.
  \item Level: an equal unit-duration long at every tenor (2y, 5y, 10y, 30y).
  \item Scale back by 5 years of duration.
\end{enumerate}

\begin{center}\small
\begin{tabularx}{\linewidth}{@{}llX@{}}
\toprule
Symbol & Source & Meaning \\
\midrule
$c_t$ & \code{FRED:DFF} & daily cash rate: FRED:DFF\_t / 100 / 252, lagged one day because the overnight rate earned on day t is set at t-1 \\
$\text{MOF\_JP:JGB\_2Y}$ & \code{MOF\_JP:JGB\_2Y} & JP\_JGB par yield at 2y \\
$\text{MOF\_JP:JGB\_5Y}$ & \code{MOF\_JP:JGB\_5Y} & JP\_JGB par yield at 5y \\
$\text{MOF\_JP:JGB\_10Y}$ & \code{MOF\_JP:JGB\_10Y} & JP\_JGB par yield at 10y \\
$\text{MOF\_JP:JGB\_30Y}$ & \code{MOF\_JP:JGB\_30Y} & JP\_JGB par yield at 30y \\
\bottomrule
\end{tabularx}
\end{center}

\noindent\textbf{Orthogonalisation}

\eqn{\tilde{f}_t = f_t - \hat{a}_\tau - \hat{b}_{1,\tau}\,\tilde{f}^{\text{rt\_us\_level}}_t}

\eqn{(\hat{a}_\tau, \hat{b}_\tau) = \arg\min_{a,b} \sum_{s=\tau-504}^{\tau-1} \Bigl(f_s - a - \sum_j b_j\,\tilde{f}^{g_j}_s\Bigr)^2}

with $g_1$ = \code{rt\_us\_level}, each taken as its own \emph{orthogonalised} series, and $\tau$ the most recent refit at or before $t$.

\subsubsection{\code{rt\_uk} --- UK Gilts}
Hierarchy level 0. Method \code{single}.

\begin{quote}\small No daily gilt curve exists anywhere in the warehouse, so this is an ETF total return converted to base currency rather than a curve factor.\end{quote}

\noindent\textbf{Construction}

\eqn{f_t = \bigl(x^{\text{IGLT.L}}_t + \Delta \ln e^{\text{GBP}}_t\bigr)}

\begin{enumerate}[leftmargin=*,itemsep=2pt]
  \item Subtract the overnight rate: x\_t = r\_t - c\_(t-1). IGLT.L is a funded long position, so its return has to be measured over cash.
  \item Convert to the base currency by adding the log change in GBP/USD. In logs the conversion is a sum, with no cross term.
\end{enumerate}

\begin{center}\small
\begin{tabularx}{\linewidth}{@{}llX@{}}
\toprule
Symbol & Source & Meaning \\
\midrule
$r^{\text{IGLT.L}}_t$ & \code{IGLT.L} & instrument return: log total return of IGLT.L \\
$c_t$ & \code{FRED:DFF} & daily cash rate: FRED:DFF\_t / 100 / 252, lagged one day because the overnight rate earned on day t is set at t-1 \\
$e^{\text{GBP}}_t$ & \code{GBP} & USD per unit of GBP; its log change converts the local return into the base currency, which is exact in logs \\
\bottomrule
\end{tabularx}
\end{center}

\noindent\textbf{Orthogonalisation}

\eqn{\tilde{f}_t = f_t - \hat{a}_\tau - \hat{b}_{1,\tau}\,\tilde{f}^{\text{rt\_us\_level}}_t}

\eqn{(\hat{a}_\tau, \hat{b}_\tau) = \arg\min_{a,b} \sum_{s=\tau-504}^{\tau-1} \Bigl(f_s - a - \sum_j b_j\,\tilde{f}^{g_j}_s\Bigr)^2}

with $g_1$ = \code{rt\_us\_level}, each taken as its own \emph{orthogonalised} series, and $\tau$ the most recent refit at or before $t$.

\subsubsection{\code{rt\_breakeven} --- Inflation Breakeven}
Hierarchy level 0. Method \code{synthetic\_bond}.

\begin{quote}\small 10y breakeven change, duration-scaled. Carry excluded: a breakeven is a spread between two yields and has no coupon of its own.\end{quote}

\noindent\textbf{Construction}

\eqn{f_t = -D_{t-1}\,\Delta y_t + \tfrac{1}{2} C_{t-1}\,(\Delta y_t)^2}

\begin{enumerate}[leftmargin=*,itemsep=2pt]
  \item Treat the series as the yield of a 10y par bond and price a day's move on it, with duration and convexity taken at t-1 so nothing about today's move leaks into its own pricing.
  \item No carry leg. A breakeven is the difference between two yields and earns no coupon of its own, so the builder adds the carry term and subtracts it again; the two cancel exactly and the factor is a pure duration-scaled change.
\end{enumerate}

\begin{center}\small
\begin{tabularx}{\linewidth}{@{}llX@{}}
\toprule
Symbol & Source & Meaning \\
\midrule
$\text{FRED:T10YIE}$ & \code{FRED:T10YIE} & quoted yield series, in percent; y\_t is that divided by 100 \\
\bottomrule
\end{tabularx}
\end{center}

\noindent\textbf{Orthogonalisation}

\eqn{\tilde{f}_t = f_t - \hat{a}_\tau - \hat{b}_{1,\tau}\,\tilde{f}^{\text{rt\_us\_level}}_t}

\eqn{(\hat{a}_\tau, \hat{b}_\tau) = \arg\min_{a,b} \sum_{s=\tau-504}^{\tau-1} \Bigl(f_s - a - \sum_j b_j\,\tilde{f}^{g_j}_s\Bigr)^2}

with $g_1$ = \code{rt\_us\_level}, each taken as its own \emph{orthogonalised} series, and $\tau$ the most recent refit at or before $t$.


% ---------------------------------------------------
\subsection{Credit}

\subsubsection{\code{cr\_ig} --- IG Credit}
Hierarchy level 1. Method \code{single}.

\begin{quote}\small Rates-hedged by regression, not by analytic duration match.\end{quote}

\noindent\textbf{Construction}

\eqn{f_t = x^{\text{LQD}}_t}

\begin{enumerate}[leftmargin=*,itemsep=2pt]
  \item Subtract the overnight rate: x\_t = r\_t - c\_(t-1). LQD is a funded long position, so its return has to be measured over cash.
\end{enumerate}

\begin{center}\small
\begin{tabularx}{\linewidth}{@{}llX@{}}
\toprule
Symbol & Source & Meaning \\
\midrule
$r^{\text{LQD}}_t$ & \code{LQD} & instrument return: log total return of LQD \\
$c_t$ & \code{FRED:DFF} & daily cash rate: FRED:DFF\_t / 100 / 252, lagged one day because the overnight rate earned on day t is set at t-1 \\
\bottomrule
\end{tabularx}
\end{center}

\noindent\textbf{Orthogonalisation}

\eqn{\tilde{f}_t = f_t - \hat{a}_\tau - \hat{b}_{1,\tau}\,\tilde{f}^{\text{rt\_us\_level}}_t - \hat{b}_{2,\tau}\,\tilde{f}^{\text{rt\_us\_slope}}_t}

\eqn{(\hat{a}_\tau, \hat{b}_\tau) = \arg\min_{a,b} \sum_{s=\tau-504}^{\tau-1} \Bigl(f_s - a - \sum_j b_j\,\tilde{f}^{g_j}_s\Bigr)^2}

with $g_1, g_2$ = \code{rt\_us\_level}, \code{rt\_us\_slope}, each taken as its own \emph{orthogonalised} series, and $\tau$ the most recent refit at or before $t$.

\subsubsection{\code{cr\_hy} --- High Yield}
Hierarchy level 1. Method \code{single}.

\begin{quote}\small Also residualised against equity: HY beta to equity is large and leaving it in would double-count equity risk.\end{quote}

\noindent\textbf{Construction}

\eqn{f_t = x^{\text{HYG}}_t}

\begin{enumerate}[leftmargin=*,itemsep=2pt]
  \item Subtract the overnight rate: x\_t = r\_t - c\_(t-1). HYG is a funded long position, so its return has to be measured over cash.
\end{enumerate}

\begin{center}\small
\begin{tabularx}{\linewidth}{@{}llX@{}}
\toprule
Symbol & Source & Meaning \\
\midrule
$r^{\text{HYG}}_t$ & \code{HYG} & instrument return: log total return of HYG \\
$c_t$ & \code{FRED:DFF} & daily cash rate: FRED:DFF\_t / 100 / 252, lagged one day because the overnight rate earned on day t is set at t-1 \\
\bottomrule
\end{tabularx}
\end{center}

\noindent\textbf{Orthogonalisation}

\eqn{\tilde{f}_t = f_t - \hat{a}_\tau - \hat{b}_{1,\tau}\,\tilde{f}^{\text{rt\_us\_level}}_t - \hat{b}_{2,\tau}\,\tilde{f}^{\text{rt\_us\_slope}}_t - \hat{b}_{3,\tau}\,\tilde{f}^{\text{eq\_global}}_t}

\eqn{(\hat{a}_\tau, \hat{b}_\tau) = \arg\min_{a,b} \sum_{s=\tau-504}^{\tau-1} \Bigl(f_s - a - \sum_j b_j\,\tilde{f}^{g_j}_s\Bigr)^2}

with $g_1, \dots, g_{3}$ = \code{rt\_us\_level}, \code{rt\_us\_slope}, \code{eq\_global}, each taken as its own \emph{orthogonalised} series, and $\tau$ the most recent refit at or before $t$.

\subsubsection{\code{cr\_loans} --- Leveraged Loans}
Hierarchy level 1. Method \code{single}.

\begin{quote}\small Floating rate, so little duration; the interesting part is what it adds over HY.\end{quote}

\noindent\textbf{Construction}

\eqn{f_t = x^{\text{BKLN}}_t}

\begin{enumerate}[leftmargin=*,itemsep=2pt]
  \item Subtract the overnight rate: x\_t = r\_t - c\_(t-1). BKLN is a funded long position, so its return has to be measured over cash.
\end{enumerate}

\begin{center}\small
\begin{tabularx}{\linewidth}{@{}llX@{}}
\toprule
Symbol & Source & Meaning \\
\midrule
$r^{\text{BKLN}}_t$ & \code{BKLN} & instrument return: log total return of BKLN \\
$c_t$ & \code{FRED:DFF} & daily cash rate: FRED:DFF\_t / 100 / 252, lagged one day because the overnight rate earned on day t is set at t-1 \\
\bottomrule
\end{tabularx}
\end{center}

\noindent\textbf{Orthogonalisation}

\eqn{\tilde{f}_t = f_t - \hat{a}_\tau - \hat{b}_{1,\tau}\,\tilde{f}^{\text{rt\_us\_level}}_t - \hat{b}_{2,\tau}\,\tilde{f}^{\text{cr\_hy}}_t}

\eqn{(\hat{a}_\tau, \hat{b}_\tau) = \arg\min_{a,b} \sum_{s=\tau-504}^{\tau-1} \Bigl(f_s - a - \sum_j b_j\,\tilde{f}^{g_j}_s\Bigr)^2}

with $g_1, g_2$ = \code{rt\_us\_level}, \code{cr\_hy}, each taken as its own \emph{orthogonalised} series, and $\tau$ the most recent refit at or before $t$.

\subsubsection{\code{cr\_em\_hard} --- EM Hard Currency}
Hierarchy level 1. Method \code{single}.

\begin{quote}\small Residualised against equity as well: EM sovereign spreads carry real equity beta (0.68 before this was removed), which would otherwise be double-counted against the equity block.\end{quote}

\noindent\textbf{Construction}

\eqn{f_t = x^{\text{EMB}}_t}

\begin{enumerate}[leftmargin=*,itemsep=2pt]
  \item Subtract the overnight rate: x\_t = r\_t - c\_(t-1). EMB is a funded long position, so its return has to be measured over cash.
\end{enumerate}

\begin{center}\small
\begin{tabularx}{\linewidth}{@{}llX@{}}
\toprule
Symbol & Source & Meaning \\
\midrule
$r^{\text{EMB}}_t$ & \code{EMB} & instrument return: log total return of EMB \\
$c_t$ & \code{FRED:DFF} & daily cash rate: FRED:DFF\_t / 100 / 252, lagged one day because the overnight rate earned on day t is set at t-1 \\
\bottomrule
\end{tabularx}
\end{center}

\noindent\textbf{Orthogonalisation}

\eqn{\tilde{f}_t = f_t - \hat{a}_\tau - \hat{b}_{1,\tau}\,\tilde{f}^{\text{rt\_us\_level}}_t - \hat{b}_{2,\tau}\,\tilde{f}^{\text{rt\_us\_slope}}_t - \hat{b}_{3,\tau}\,\tilde{f}^{\text{cr\_hy}}_t - \hat{b}_{4,\tau}\,\tilde{f}^{\text{eq\_global}}_t}

\eqn{(\hat{a}_\tau, \hat{b}_\tau) = \arg\min_{a,b} \sum_{s=\tau-504}^{\tau-1} \Bigl(f_s - a - \sum_j b_j\,\tilde{f}^{g_j}_s\Bigr)^2}

with $g_1, \dots, g_{4}$ = \code{rt\_us\_level}, \code{rt\_us\_slope}, \code{cr\_hy}, \code{eq\_global}, each taken as its own \emph{orthogonalised} series, and $\tau$ the most recent refit at or before $t$.

\subsubsection{\code{cr\_em\_local} --- EM Local Currency}
Hierarchy level 2. Method \code{single}.

\begin{quote}\small Residualised against USD as well: the local-currency leg is mostly FX.\end{quote}

\noindent\textbf{Construction}

\eqn{f_t = x^{\text{EMLC}}_t}

\begin{enumerate}[leftmargin=*,itemsep=2pt]
  \item Subtract the overnight rate: x\_t = r\_t - c\_(t-1). EMLC is a funded long position, so its return has to be measured over cash.
\end{enumerate}

\begin{center}\small
\begin{tabularx}{\linewidth}{@{}llX@{}}
\toprule
Symbol & Source & Meaning \\
\midrule
$r^{\text{EMLC}}_t$ & \code{EMLC} & instrument return: log total return of EMLC \\
$c_t$ & \code{FRED:DFF} & daily cash rate: FRED:DFF\_t / 100 / 252, lagged one day because the overnight rate earned on day t is set at t-1 \\
\bottomrule
\end{tabularx}
\end{center}

\noindent\textbf{Orthogonalisation}

\eqn{\tilde{f}_t = f_t - \hat{a}_\tau - \hat{b}_{1,\tau}\,\tilde{f}^{\text{rt\_us\_level}}_t - \hat{b}_{2,\tau}\,\tilde{f}^{\text{cr\_em\_hard}}_t - \hat{b}_{3,\tau}\,\tilde{f}^{\text{fx\_usd}}_t}

\eqn{(\hat{a}_\tau, \hat{b}_\tau) = \arg\min_{a,b} \sum_{s=\tau-504}^{\tau-1} \Bigl(f_s - a - \sum_j b_j\,\tilde{f}^{g_j}_s\Bigr)^2}

with $g_1, \dots, g_{3}$ = \code{rt\_us\_level}, \code{cr\_em\_hard}, \code{fx\_usd}, each taken as its own \emph{orthogonalised} series, and $\tau$ the most recent refit at or before $t$.


% ---------------------------------------------------
\subsection{FX}

\subsubsection{\code{fx\_usd} --- USD Broad}
Hierarchy level 0. Method \code{single}.

\begin{quote}\small Dollar index return. Not a funded position, so no cash leg.\end{quote}

\noindent\textbf{Construction}

\eqn{f_t = r^{\text{DX-Y.NYB}}_t}

\begin{enumerate}[leftmargin=*,itemsep=2pt]
  \item No cash leg: this is not a funded position, so there is nothing to fund.
\end{enumerate}

\begin{center}\small
\begin{tabularx}{\linewidth}{@{}llX@{}}
\toprule
Symbol & Source & Meaning \\
\midrule
$r^{\text{DX-Y.NYB}}_t$ & \code{DX-Y.NYB} & instrument return: log total return of DX-Y.NYB \\
\bottomrule
\end{tabularx}
\end{center}

\noindent\textbf{Orthogonalisation}

\eqn{\tilde{f}_t = f_t}

Nothing sits above this factor in the hierarchy, so the raw and orthogonalised panels hold the same series.

\subsubsection{\code{fx\_jpy} --- JPY (ex-USD)}
Hierarchy level 1. Method \code{single}.

\begin{quote}\small Sign flipped so a positive value means a stronger yen, the usual risk-off direction.\end{quote}

\noindent\textbf{Construction}

\eqn{f_t = -1\,r^{\text{USDJPY=X}}_t}

\begin{enumerate}[leftmargin=*,itemsep=2pt]
  \item No cash leg: this is not a funded position, so there is nothing to fund.
  \item Multiply by -1 so the factor points in the stated direction.
\end{enumerate}

\begin{center}\small
\begin{tabularx}{\linewidth}{@{}llX@{}}
\toprule
Symbol & Source & Meaning \\
\midrule
$r^{\text{USDJPY=X}}_t$ & \code{USDJPY=X} & instrument return: log total return of USDJPY=X \\
\bottomrule
\end{tabularx}
\end{center}

\noindent\textbf{Orthogonalisation}

\eqn{\tilde{f}_t = f_t - \hat{a}_\tau - \hat{b}_{1,\tau}\,\tilde{f}^{\text{fx\_usd}}_t}

\eqn{(\hat{a}_\tau, \hat{b}_\tau) = \arg\min_{a,b} \sum_{s=\tau-504}^{\tau-1} \Bigl(f_s - a - \sum_j b_j\,\tilde{f}^{g_j}_s\Bigr)^2}

with $g_1$ = \code{fx\_usd}, each taken as its own \emph{orthogonalised} series, and $\tau$ the most recent refit at or before $t$.

\subsubsection{\code{fx\_em} --- EM FX}
Hierarchy level 1. Method \code{single}.

\noindent\textbf{Construction}

\eqn{f_t = x^{\text{CEW}}_t}

\begin{enumerate}[leftmargin=*,itemsep=2pt]
  \item Subtract the overnight rate: x\_t = r\_t - c\_(t-1). CEW is a funded long position, so its return has to be measured over cash.
\end{enumerate}

\begin{center}\small
\begin{tabularx}{\linewidth}{@{}llX@{}}
\toprule
Symbol & Source & Meaning \\
\midrule
$r^{\text{CEW}}_t$ & \code{CEW} & instrument return: log total return of CEW \\
$c_t$ & \code{FRED:DFF} & daily cash rate: FRED:DFF\_t / 100 / 252, lagged one day because the overnight rate earned on day t is set at t-1 \\
\bottomrule
\end{tabularx}
\end{center}

\noindent\textbf{Orthogonalisation}

\eqn{\tilde{f}_t = f_t - \hat{a}_\tau - \hat{b}_{1,\tau}\,\tilde{f}^{\text{fx\_usd}}_t}

\eqn{(\hat{a}_\tau, \hat{b}_\tau) = \arg\min_{a,b} \sum_{s=\tau-504}^{\tau-1} \Bigl(f_s - a - \sum_j b_j\,\tilde{f}^{g_j}_s\Bigr)^2}

with $g_1$ = \code{fx\_usd}, each taken as its own \emph{orthogonalised} series, and $\tau$ the most recent refit at or before $t$.

\subsubsection{\code{fx\_carry} --- FX Carry (G4 approximation)}
Hierarchy level 1. Method \code{fx\_carry}.

\begin{quote}\small APPROXIMATION. Built from policy-rate differentials because the warehouse has no forward points, and limited to the three non-USD policy rates available. Narrower and noisier than a real G10 carry basket; treat its loading with corresponding scepticism. Also residualised against JPY: with only three legs and the yen as the perennial funding currency, the raw basket correlated -0.87 with fx\_jpy, i.e. it was the yen short wearing a different name.\end{quote}

\noindent\textbf{Construction}

\eqn{\kappa_{j,t} = \frac{i_{j,t-1} - i_{\text{USD},t-1}}{100 \cdot 252}, \qquad p_{j,t} = \operatorname{sign}(\kappa_{j,t})}

\eqn{f_t = \frac{1}{3}\sum_j p_{j,t}\bigl(s_j\,r^{X_j}_t + \kappa_{j,t}\bigr)}

\begin{enumerate}[leftmargin=*,itemsep=2pt]
  \item Carry is approximated from policy-rate differentials (CHF, EUR, JPY against USD) because the warehouse holds no forward points. Covered interest parity says the forward discount equals the rate differential, so this stands in for it.
  \item The differential is lagged one day: the carry earned over day t is fixed by the rates set at t-1.
  \item Each leg is held long when it yields more than the dollar and short when it yields less, and the basket is the equally weighted average of the legs.
  \item APPROXIMATION. Three non-USD legs is narrower and noisier than a real G10 carry basket; treat its loading with corresponding scepticism.
\end{enumerate}

\begin{center}\small
\begin{tabularx}{\linewidth}{@{}llX@{}}
\toprule
Symbol & Source & Meaning \\
\midrule
$r^{\text{USDCHF=X}}_t$ & \code{USDCHF=X} & CHF spot leg, sign flipped so the quoted pair reads as long the foreign currency against the dollar: log total return of USDCHF=X \\
$\text{SNB:POLICY\_RATE}$ & \code{SNB:POLICY\_RATE} & CHF policy rate i\_CHF, annualised percent \\
$r^{\text{EURUSD=X}}_t$ & \code{EURUSD=X} & EUR spot leg: log total return of EURUSD=X \\
$\text{ECB:DFR}$ & \code{ECB:DFR} & EUR policy rate i\_EUR, annualised percent \\
$r^{\text{USDJPY=X}}_t$ & \code{USDJPY=X} & JPY spot leg, sign flipped so the quoted pair reads as long the foreign currency against the dollar: log total return of USDJPY=X \\
$\text{BOJ:IR01\_OCRT}$ & \code{BOJ:IR01\_OCRT} & JPY policy rate i\_JPY, annualised percent \\
$\text{FRED:DFF}$ & \code{FRED:DFF} & USD policy rate, the funding leg \\
\bottomrule
\end{tabularx}
\end{center}

\noindent\textbf{Orthogonalisation}

\eqn{\tilde{f}_t = f_t - \hat{a}_\tau - \hat{b}_{1,\tau}\,\tilde{f}^{\text{fx\_usd}}_t - \hat{b}_{2,\tau}\,\tilde{f}^{\text{fx\_jpy}}_t}

\eqn{(\hat{a}_\tau, \hat{b}_\tau) = \arg\min_{a,b} \sum_{s=\tau-504}^{\tau-1} \Bigl(f_s - a - \sum_j b_j\,\tilde{f}^{g_j}_s\Bigr)^2}

with $g_1, g_2$ = \code{fx\_usd}, \code{fx\_jpy}, each taken as its own \emph{orthogonalised} series, and $\tau$ the most recent refit at or before $t$.


% ---------------------------------------------------
\subsection{Commodities}

\subsubsection{\code{cm\_broad} --- Broad Commodity}
Hierarchy level 0. Method \code{single}.

\begin{quote}\small DBC is a total-return futures fund: roll and collateral included.\end{quote}

\noindent\textbf{Construction}

\eqn{f_t = x^{\text{DBC}}_t}

\begin{enumerate}[leftmargin=*,itemsep=2pt]
  \item Subtract the overnight rate: x\_t = r\_t - c\_(t-1). DBC is a funded long position, so its return has to be measured over cash.
\end{enumerate}

\begin{center}\small
\begin{tabularx}{\linewidth}{@{}llX@{}}
\toprule
Symbol & Source & Meaning \\
\midrule
$r^{\text{DBC}}_t$ & \code{DBC} & instrument return: log total return of DBC \\
$c_t$ & \code{FRED:DFF} & daily cash rate: FRED:DFF\_t / 100 / 252, lagged one day because the overnight rate earned on day t is set at t-1 \\
\bottomrule
\end{tabularx}
\end{center}

\noindent\textbf{Orthogonalisation}

\eqn{\tilde{f}_t = f_t}

Nothing sits above this factor in the hierarchy, so the raw and orthogonalised panels hold the same series.

\subsubsection{\code{cm\_energy} --- Energy (ex-broad)}
Hierarchy level 1. Method \code{single}.

\noindent\textbf{Construction}

\eqn{f_t = x^{\text{USO}}_t}

\begin{enumerate}[leftmargin=*,itemsep=2pt]
  \item Subtract the overnight rate: x\_t = r\_t - c\_(t-1). USO is a funded long position, so its return has to be measured over cash.
\end{enumerate}

\begin{center}\small
\begin{tabularx}{\linewidth}{@{}llX@{}}
\toprule
Symbol & Source & Meaning \\
\midrule
$r^{\text{USO}}_t$ & \code{USO} & instrument return: log total return of USO \\
$c_t$ & \code{FRED:DFF} & daily cash rate: FRED:DFF\_t / 100 / 252, lagged one day because the overnight rate earned on day t is set at t-1 \\
\bottomrule
\end{tabularx}
\end{center}

\noindent\textbf{Orthogonalisation}

\eqn{\tilde{f}_t = f_t - \hat{a}_\tau - \hat{b}_{1,\tau}\,\tilde{f}^{\text{cm\_broad}}_t}

\eqn{(\hat{a}_\tau, \hat{b}_\tau) = \arg\min_{a,b} \sum_{s=\tau-504}^{\tau-1} \Bigl(f_s - a - \sum_j b_j\,\tilde{f}^{g_j}_s\Bigr)^2}

with $g_1$ = \code{cm\_broad}, each taken as its own \emph{orthogonalised} series, and $\tau$ the most recent refit at or before $t$.

\subsubsection{\code{cm\_gold} --- Gold (ex-broad)}
Hierarchy level 1. Method \code{single}.

\noindent\textbf{Construction}

\eqn{f_t = x^{\text{GLD}}_t}

\begin{enumerate}[leftmargin=*,itemsep=2pt]
  \item Subtract the overnight rate: x\_t = r\_t - c\_(t-1). GLD is a funded long position, so its return has to be measured over cash.
\end{enumerate}

\begin{center}\small
\begin{tabularx}{\linewidth}{@{}llX@{}}
\toprule
Symbol & Source & Meaning \\
\midrule
$r^{\text{GLD}}_t$ & \code{GLD} & instrument return: log total return of GLD \\
$c_t$ & \code{FRED:DFF} & daily cash rate: FRED:DFF\_t / 100 / 252, lagged one day because the overnight rate earned on day t is set at t-1 \\
\bottomrule
\end{tabularx}
\end{center}

\noindent\textbf{Orthogonalisation}

\eqn{\tilde{f}_t = f_t - \hat{a}_\tau - \hat{b}_{1,\tau}\,\tilde{f}^{\text{cm\_broad}}_t}

\eqn{(\hat{a}_\tau, \hat{b}_\tau) = \arg\min_{a,b} \sum_{s=\tau-504}^{\tau-1} \Bigl(f_s - a - \sum_j b_j\,\tilde{f}^{g_j}_s\Bigr)^2}

with $g_1$ = \code{cm\_broad}, each taken as its own \emph{orthogonalised} series, and $\tau$ the most recent refit at or before $t$.

\subsubsection{\code{cm\_industrial} --- Industrial Metals (ex-broad)}
Hierarchy level 1. Method \code{single}.

\noindent\textbf{Construction}

\eqn{f_t = x^{\text{CPER}}_t}

\begin{enumerate}[leftmargin=*,itemsep=2pt]
  \item Subtract the overnight rate: x\_t = r\_t - c\_(t-1). CPER is a funded long position, so its return has to be measured over cash.
\end{enumerate}

\begin{center}\small
\begin{tabularx}{\linewidth}{@{}llX@{}}
\toprule
Symbol & Source & Meaning \\
\midrule
$r^{\text{CPER}}_t$ & \code{CPER} & instrument return: log total return of CPER \\
$c_t$ & \code{FRED:DFF} & daily cash rate: FRED:DFF\_t / 100 / 252, lagged one day because the overnight rate earned on day t is set at t-1 \\
\bottomrule
\end{tabularx}
\end{center}

\noindent\textbf{Orthogonalisation}

\eqn{\tilde{f}_t = f_t - \hat{a}_\tau - \hat{b}_{1,\tau}\,\tilde{f}^{\text{cm\_broad}}_t}

\eqn{(\hat{a}_\tau, \hat{b}_\tau) = \arg\min_{a,b} \sum_{s=\tau-504}^{\tau-1} \Bigl(f_s - a - \sum_j b_j\,\tilde{f}^{g_j}_s\Bigr)^2}

with $g_1$ = \code{cm\_broad}, each taken as its own \emph{orthogonalised} series, and $\tau$ the most recent refit at or before $t$.

\subsubsection{\code{cm\_agri} --- Agriculture (ex-broad)}
Hierarchy level 1. Method \code{single}.

\noindent\textbf{Construction}

\eqn{f_t = x^{\text{DBA}}_t}

\begin{enumerate}[leftmargin=*,itemsep=2pt]
  \item Subtract the overnight rate: x\_t = r\_t - c\_(t-1). DBA is a funded long position, so its return has to be measured over cash.
\end{enumerate}

\begin{center}\small
\begin{tabularx}{\linewidth}{@{}llX@{}}
\toprule
Symbol & Source & Meaning \\
\midrule
$r^{\text{DBA}}_t$ & \code{DBA} & instrument return: log total return of DBA \\
$c_t$ & \code{FRED:DFF} & daily cash rate: FRED:DFF\_t / 100 / 252, lagged one day because the overnight rate earned on day t is set at t-1 \\
\bottomrule
\end{tabularx}
\end{center}

\noindent\textbf{Orthogonalisation}

\eqn{\tilde{f}_t = f_t - \hat{a}_\tau - \hat{b}_{1,\tau}\,\tilde{f}^{\text{cm\_broad}}_t}

\eqn{(\hat{a}_\tau, \hat{b}_\tau) = \arg\min_{a,b} \sum_{s=\tau-504}^{\tau-1} \Bigl(f_s - a - \sum_j b_j\,\tilde{f}^{g_j}_s\Bigr)^2}

with $g_1$ = \code{cm\_broad}, each taken as its own \emph{orthogonalised} series, and $\tau$ the most recent refit at or before $t$.


% ---------------------------------------------------
\subsection{Volatility}

\subsubsection{\code{vol\_equity} --- Equity Volatility (VIX futures)}
Hierarchy level 1. Method \code{single}.

\begin{quote}\small Long VIX-futures strategy return, residualised against equity.\end{quote}

\noindent\textbf{Construction}

\eqn{f_t = x^{\text{VIXY}}_t}

\begin{enumerate}[leftmargin=*,itemsep=2pt]
  \item Subtract the overnight rate: x\_t = r\_t - c\_(t-1). VIXY is a funded long position, so its return has to be measured over cash.
\end{enumerate}

\begin{center}\small
\begin{tabularx}{\linewidth}{@{}llX@{}}
\toprule
Symbol & Source & Meaning \\
\midrule
$r^{\text{VIXY}}_t$ & \code{VIXY} & instrument return: log total return of VIXY \\
$c_t$ & \code{FRED:DFF} & daily cash rate: FRED:DFF\_t / 100 / 252, lagged one day because the overnight rate earned on day t is set at t-1 \\
\bottomrule
\end{tabularx}
\end{center}

\noindent\textbf{Orthogonalisation}

\eqn{\tilde{f}_t = f_t - \hat{a}_\tau - \hat{b}_{1,\tau}\,\tilde{f}^{\text{eq\_global}}_t}

\eqn{(\hat{a}_\tau, \hat{b}_\tau) = \arg\min_{a,b} \sum_{s=\tau-504}^{\tau-1} \Bigl(f_s - a - \sum_j b_j\,\tilde{f}^{g_j}_s\Bigr)^2}

with $g_1$ = \code{eq\_global}, each taken as its own \emph{orthogonalised} series, and $\tau$ the most recent refit at or before $t$.

\subsubsection{\code{vol\_variance\_premium} --- Variance Risk Premium}
Hierarchy level 1. Method \code{variance\_premium}.

\begin{quote}\small Daily payoff of a short variance position: yesterday's implied variance minus today's realised squared return.\end{quote}

\noindent\textbf{Construction}

\eqn{f_t = \frac{1}{252}\left(\frac{\text{VIX}_{t-1}}{100}\right)^2 - \bigl(r^{\text{SPY}}_t\bigr)^2}

\eqn{f_t \leftarrow k\,f_t, \qquad k = \frac{0.1 / \sqrt{252}}{\operatorname{sd}(f)}}

\begin{enumerate}[leftmargin=*,itemsep=2pt]
  \item Yesterday's implied variance is what a variance swap struck at t-1 pays against, so the lag is the contract and not a modelling choice.
  \item Today's realised variance is the squared SPY return, which is what the daily leg of such a swap settles on.
  \item The difference is the daily payoff of a short variance position, positive on average because implied sits above realised.
  \item The payoff is convex in the underlying, so a linear beta on it is a first approximation and nothing more.
  \item Rescale to 10\% annualised volatility. The transform above yields a z-score, and every factor in this model is a return.
\end{enumerate}

\begin{center}\small
\begin{tabularx}{\linewidth}{@{}llX@{}}
\toprule
Symbol & Source & Meaning \\
\midrule
$\text{FRED:VIXCLS}$ & \code{FRED:VIXCLS} & VIX close in index points; divided by 100 it is an annualised volatility in decimal \\
$r^{\text{SPY}}_t$ & \code{SPY} & underlying whose realised variance is delivered: log total return of SPY \\
\bottomrule
\end{tabularx}
\end{center}

\noindent\textbf{Orthogonalisation}

\eqn{\tilde{f}_t = f_t - \hat{a}_\tau - \hat{b}_{1,\tau}\,\tilde{f}^{\text{eq\_global}}_t}

\eqn{(\hat{a}_\tau, \hat{b}_\tau) = \arg\min_{a,b} \sum_{s=\tau-504}^{\tau-1} \Bigl(f_s - a - \sum_j b_j\,\tilde{f}^{g_j}_s\Bigr)^2}

with $g_1$ = \code{eq\_global}, each taken as its own \emph{orthogonalised} series, and $\tau$ the most recent refit at or before $t$.

\subsubsection{\code{vol\_rates} --- Rates Volatility (realised)}
Hierarchy level 1. Method \code{realized\_vol\_change}.

\begin{quote}\small Substitute for \textasciicircum{}TYVIX, which CBOE discontinued in 2020.\end{quote}

\noindent\textbf{Construction}

\eqn{v_t = \sqrt{252}\;\operatorname{sd}\bigl(r^{\text{TLT}}_{t-20},\dots,r^{\text{TLT}}_t\bigr)}

\eqn{f_t = \frac{\Delta v_t}{\operatorname{sd}(\Delta v_{t-251},\dots,\Delta v_t)}}

\eqn{f_t \leftarrow k\,f_t, \qquad k = \frac{0.1 / \sqrt{252}}{\operatorname{sd}(f)}}

\begin{enumerate}[leftmargin=*,itemsep=2pt]
  \item Realised volatility over a trailing 21 days, annualised.
  \item Differenced, because a volatility level is not stationary.
  \item Standardised on a trailing window, because a change in volatility is not measured in return units.
  \item This is a change in a risk measure, not a tradable return. It stands in for a rates-volatility series because CBOE discontinued \textasciicircum{}TYVIX in 2020.
  \item Rescale to 10\% annualised volatility. The transform above yields a z-score, and every factor in this model is a return.
\end{enumerate}

\begin{center}\small
\begin{tabularx}{\linewidth}{@{}llX@{}}
\toprule
Symbol & Source & Meaning \\
\midrule
$r^{\text{TLT}}_t$ & \code{TLT} & instrument whose realised volatility is tracked: log total return of TLT \\
\bottomrule
\end{tabularx}
\end{center}

\noindent\textbf{Orthogonalisation}

\eqn{\tilde{f}_t = f_t - \hat{a}_\tau - \hat{b}_{1,\tau}\,\tilde{f}^{\text{rt\_us\_level}}_t}

\eqn{(\hat{a}_\tau, \hat{b}_\tau) = \arg\min_{a,b} \sum_{s=\tau-504}^{\tau-1} \Bigl(f_s - a - \sum_j b_j\,\tilde{f}^{g_j}_s\Bigr)^2}

with $g_1$ = \code{rt\_us\_level}, each taken as its own \emph{orthogonalised} series, and $\tau$ the most recent refit at or before $t$.


% ---------------------------------------------------
\subsection{Liquidity \& Stress}

\subsubsection{\code{liq\_funding} --- Funding Spread (SOFR-EFFR)}
Hierarchy level 0. Method \code{spread\_level}.

\begin{quote}\small TED and LIBOR-OIS are discontinued; SOFR over the effective funds rate is the live equivalent. Standardised because basis points are not comparable to the return-scaled factors.\end{quote}

\noindent\textbf{Construction}

\eqn{s_t = \text{FRED:SOFR}_t - \text{FRED:EFFR}_t}

\eqn{f_t = \frac{\Delta s_t}{\operatorname{sd}\bigl(\Delta s_{t-251},\dots,\Delta s_t\bigr)}}

\eqn{f_t \leftarrow k\,f_t, \qquad k = \frac{0.1 / \sqrt{252}}{\operatorname{sd}(f)}}

\begin{enumerate}[leftmargin=*,itemsep=2pt]
  \item Form the spread between the two quoted levels.
  \item Divide the daily change by its own trailing volatility. Basis points are not comparable with the return-scaled blocks, and the scaling window is strictly trailing: a full-sample standard deviation would leak future volatility into every historical observation.
  \item Rescale to 10\% annualised volatility. The transform above yields a z-score, and every factor in this model is a return.
\end{enumerate}

\begin{center}\small
\begin{tabularx}{\linewidth}{@{}llX@{}}
\toprule
Symbol & Source & Meaning \\
\midrule
$\text{FRED:SOFR}$ & \code{FRED:SOFR} & first leg of the spread, in percent \\
$\text{FRED:EFFR}$ & \code{FRED:EFFR} & second leg of the spread, in percent \\
\bottomrule
\end{tabularx}
\end{center}

\noindent\textbf{Orthogonalisation}

\eqn{\tilde{f}_t = f_t}

Nothing sits above this factor in the hierarchy, so the raw and orthogonalised panels hold the same series.

\subsubsection{\code{liq\_conditions} --- Financial Conditions (NFCI)}
Hierarchy level 0. Method \code{level\_transform}.

\begin{quote}\small Weekly series; the daily factor is zero between releases rather than forward-filled, since a forward fill would manufacture information.\end{quote}

\noindent\textbf{Construction}

\eqn{f_t = \begin{cases}\dfrac{\Delta s_k}{\operatorname{sd}(\Delta s_{k-51},\dots,\Delta s_k)} & t = \tau_k,\ \text{a release day} \\[2ex]0 & \text{otherwise}\end{cases}}

\eqn{f_t \leftarrow k\,f_t, \qquad k = \frac{0.1 / \sqrt{252}}{\operatorname{sd}(f)}}

\begin{enumerate}[leftmargin=*,itemsep=2pt]
  \item The series prints weekly, not daily.
  \item It is NOT forward filled. A filled series carries no new information between releases, which manufactures autocorrelation, understates standard errors and can smuggle in lookahead.
  \item Instead the standardised change lands on the publication day and the factor is exactly zero in between: a release-event factor, sparse daily.
  \item The scaling window counts releases rather than days, so 52 of them is about a year of the series' own history.
  \item Rescale to 10\% annualised volatility. The transform above yields a z-score, and every factor in this model is a return.
\end{enumerate}

\begin{center}\small
\begin{tabularx}{\linewidth}{@{}llX@{}}
\toprule
Symbol & Source & Meaning \\
\midrule
$\text{FRED:NFCI}$ & \code{FRED:NFCI} & the published level series \\
\bottomrule
\end{tabularx}
\end{center}

\noindent\textbf{Orthogonalisation}

\eqn{\tilde{f}_t = f_t}

Nothing sits above this factor in the hierarchy, so the raw and orthogonalised panels hold the same series.

\subsubsection{\code{liq\_risk\_off} --- Cross-Asset Risk-Off}
Hierarchy level 2. Method \code{basket}.

\begin{quote}\small The residual co-movement of classic haven and risk assets after each leg's own block exposure is removed.\end{quote}

\noindent\textbf{Construction}

\eqn{f_t = \frac{r^{\text{TLT}}_t + r^{\text{UUP}}_t + r^{\text{GLD}}_t}{3} - \frac{r^{\text{HYG}}_t + r^{\text{EEM}}_t}{2}}

\begin{enumerate}[leftmargin=*,itemsep=2pt]
  \item Equally weighted long basket of 3: each leg carries 1/3 of the notional.
  \item Minus an equally weighted short basket of 2. Both sides carry the market direction, so it cancels and what survives is the spread between them.
  \item Self-financing, so no cash leg.
\end{enumerate}

\begin{center}\small
\begin{tabularx}{\linewidth}{@{}llX@{}}
\toprule
Symbol & Source & Meaning \\
\midrule
$r^{\text{TLT}}_t$ & \code{TLT} & long leg: log total return of TLT \\
$r^{\text{UUP}}_t$ & \code{UUP} & long leg: log total return of UUP \\
$r^{\text{GLD}}_t$ & \code{GLD} & long leg: log total return of GLD \\
$r^{\text{HYG}}_t$ & \code{HYG} & short leg: log total return of HYG \\
$r^{\text{EEM}}_t$ & \code{EEM} & short leg: log total return of EEM \\
\bottomrule
\end{tabularx}
\end{center}

\noindent\textbf{Orthogonalisation}

\eqn{\tilde{f}_t = f_t - \hat{a}_\tau - \hat{b}_{1,\tau}\,\tilde{f}^{\text{eq\_global}}_t - \hat{b}_{2,\tau}\,\tilde{f}^{\text{rt\_us\_level}}_t - \hat{b}_{3,\tau}\,\tilde{f}^{\text{cr\_hy}}_t}

\eqn{(\hat{a}_\tau, \hat{b}_\tau) = \arg\min_{a,b} \sum_{s=\tau-504}^{\tau-1} \Bigl(f_s - a - \sum_j b_j\,\tilde{f}^{g_j}_s\Bigr)^2}

with $g_1, \dots, g_{3}$ = \code{eq\_global}, \code{rt\_us\_level}, \code{cr\_hy}, each taken as its own \emph{orthogonalised} series, and $\tau$ the most recent refit at or before $t$.


% ---------------------------------------------------
\subsection{Alternative Risk Premia}

\subsubsection{\code{arp\_trend} --- Time-Series Momentum}
Hierarchy level 3. Method \code{tsmom}.

\begin{quote}\small 12-month-minus-1 signal, inverse-volatility sized, equally weighted across assets. Validated against DBMF and KMLM over their short lives.\end{quote}

\noindent\textbf{Construction}

\eqn{S_{i,t} = \operatorname{sign}\Bigl(\sum_{u=t-252}^{t-22} r_{i,u}\Bigr), \qquad v_{i,t} = \operatorname{sd}\bigl(r_{i,t-63},\dots,r_{i,t-1}\bigr)}

\eqn{w_{i,t} = \frac{S_{i,t}}{v_{i,t}}, \qquad f_t = \frac{\sum_i w_{i,t}\,r_{i,t}}{\sum_i \lvert w_{i,t}\rvert}}

\begin{enumerate}[leftmargin=*,itemsep=2pt]
  \item The signal at t reads returns through t-22 only, and the position it implies is applied to the return at t. Nothing about day t enters its own weight.
  \item Skipping the most recent 21 days is the standard 12-minus-1 construction, which keeps short-term reversal out of the trend signal.
  \item Inverse-volatility sizing puts each of the 13 instruments on a comparable risk footing; without it the equity legs would drown out the bond legs.
  \item Dividing by gross exposure makes the factor a return on one unit of notional rather than a sum that grows with the size of the universe.
\end{enumerate}

\begin{center}\small
\begin{tabularx}{\linewidth}{@{}llX@{}}
\toprule
Symbol & Source & Meaning \\
\midrule
$r^{\text{SPY}}_t$ & \code{SPY} & instrument return: log total return of SPY \\
$r^{\text{EFA}}_t$ & \code{EFA} & instrument return: log total return of EFA \\
$r^{\text{EEM}}_t$ & \code{EEM} & instrument return: log total return of EEM \\
$r^{\text{TLT}}_t$ & \code{TLT} & instrument return: log total return of TLT \\
$r^{\text{IEF}}_t$ & \code{IEF} & instrument return: log total return of IEF \\
$r^{\text{DBC}}_t$ & \code{DBC} & instrument return: log total return of DBC \\
$r^{\text{GLD}}_t$ & \code{GLD} & instrument return: log total return of GLD \\
$r^{\text{USO}}_t$ & \code{USO} & instrument return: log total return of USO \\
$r^{\text{HYG}}_t$ & \code{HYG} & instrument return: log total return of HYG \\
$r^{\text{LQD}}_t$ & \code{LQD} & instrument return: log total return of LQD \\
$r^{\text{UUP}}_t$ & \code{UUP} & instrument return: log total return of UUP \\
$r^{\text{FXE}}_t$ & \code{FXE} & instrument return: log total return of FXE \\
$r^{\text{FXY}}_t$ & \code{FXY} & instrument return: log total return of FXY \\
\bottomrule
\end{tabularx}
\end{center}

\noindent\textbf{Orthogonalisation}

\eqn{\tilde{f}_t = f_t - \hat{a}_\tau - \hat{b}_{1,\tau}\,\tilde{f}^{\text{eq\_global}}_t - \hat{b}_{2,\tau}\,\tilde{f}^{\text{rt\_us\_level}}_t - \hat{b}_{3,\tau}\,\tilde{f}^{\text{cm\_broad}}_t - \hat{b}_{4,\tau}\,\tilde{f}^{\text{fx\_usd}}_t}

\eqn{(\hat{a}_\tau, \hat{b}_\tau) = \arg\min_{a,b} \sum_{s=\tau-504}^{\tau-1} \Bigl(f_s - a - \sum_j b_j\,\tilde{f}^{g_j}_s\Bigr)^2}

with $g_1, \dots, g_{4}$ = \code{eq\_global}, \code{rt\_us\_level}, \code{cm\_broad}, \code{fx\_usd}, each taken as its own \emph{orthogonalised} series, and $\tau$ the most recent refit at or before $t$.

\subsubsection{\code{arp\_xs\_momentum} --- Cross-Sectional Momentum}
Hierarchy level 3. Method \code{xs\_momentum}.

\begin{quote}\small Long the top third, short the bottom third by trailing return; dollar-neutral by construction.\end{quote}

\noindent\textbf{Construction}

\eqn{M_{i,t} = \sum_{u=t-252}^{t-22} r_{i,u}, \qquad q_{i,t} = \operatorname{rank}_\text{pct}\bigl(M_{\cdot,t}\bigr)_i}

\eqn{w_{i,t} = \frac{\mathbf{1}\{q_{i,t} > 2/3\}}{n^{\text{long}}_t} - \frac{\mathbf{1}\{q_{i,t} < 1/3\}}{n^{\text{short}}_t}, \qquad f_t = \sum_i w_{i,t}\,r_{i,t}}

\begin{enumerate}[leftmargin=*,itemsep=2pt]
  \item Rank the 13 instruments by trailing return, again skipping the most recent 21 days.
  \item Long the top third, short the bottom third, equally weighted within each leg.
  \item Dollar-neutral by construction: the long and the short weights each sum to one, so the common market direction cancels and what is left is the spread between winners and losers.
\end{enumerate}

\begin{center}\small
\begin{tabularx}{\linewidth}{@{}llX@{}}
\toprule
Symbol & Source & Meaning \\
\midrule
$r^{\text{SPY}}_t$ & \code{SPY} & instrument return: log total return of SPY \\
$r^{\text{EFA}}_t$ & \code{EFA} & instrument return: log total return of EFA \\
$r^{\text{EEM}}_t$ & \code{EEM} & instrument return: log total return of EEM \\
$r^{\text{TLT}}_t$ & \code{TLT} & instrument return: log total return of TLT \\
$r^{\text{IEF}}_t$ & \code{IEF} & instrument return: log total return of IEF \\
$r^{\text{DBC}}_t$ & \code{DBC} & instrument return: log total return of DBC \\
$r^{\text{GLD}}_t$ & \code{GLD} & instrument return: log total return of GLD \\
$r^{\text{USO}}_t$ & \code{USO} & instrument return: log total return of USO \\
$r^{\text{HYG}}_t$ & \code{HYG} & instrument return: log total return of HYG \\
$r^{\text{LQD}}_t$ & \code{LQD} & instrument return: log total return of LQD \\
$r^{\text{UUP}}_t$ & \code{UUP} & instrument return: log total return of UUP \\
$r^{\text{FXE}}_t$ & \code{FXE} & instrument return: log total return of FXE \\
$r^{\text{FXY}}_t$ & \code{FXY} & instrument return: log total return of FXY \\
\bottomrule
\end{tabularx}
\end{center}

\noindent\textbf{Orthogonalisation}

\eqn{\tilde{f}_t = f_t - \hat{a}_\tau - \hat{b}_{1,\tau}\,\tilde{f}^{\text{eq\_global}}_t - \hat{b}_{2,\tau}\,\tilde{f}^{\text{arp\_trend}}_t}

\eqn{(\hat{a}_\tau, \hat{b}_\tau) = \arg\min_{a,b} \sum_{s=\tau-504}^{\tau-1} \Bigl(f_s - a - \sum_j b_j\,\tilde{f}^{g_j}_s\Bigr)^2}

with $g_1, g_2$ = \code{eq\_global}, \code{arp\_trend}, each taken as its own \emph{orthogonalised} series, and $\tau$ the most recent refit at or before $t$.


% ---------------------------------------------------
\subsection{The orthogonalisation, once}

The equation above is the same for every factor that has targets; only the regressor list changes. In general form, for $k$ targets $g_1, \dots, g_k$:

\eqn{\tilde{f}_t = f_t - \hat{a}_\tau - \sum_{j=1}^{k} \hat{b}_{j,\tau}\,\tilde{f}^{g_j}_t}

\eqn{(\hat{a}_\tau, \hat{b}_\tau) = \arg\min_{a,b} \sum_{s=\tau-504}^{\tau-1} \Bigl(f_s - a - \sum_j b_j\,\tilde{f}^{g_j}_s\Bigr)^2}

\begin{enumerate}[leftmargin=*,itemsep=2pt]
  \item Coefficients come from ordinary least squares on the trailing 504 observations (2 years), refitted every 21 observations and applied forward until the next refit.
  \item The fitting window ends at tau-1, so day t is never part of the regression that residualises it. That is the whole reason for the rolling refit: a full-sample residual would put future information into historical factor values, and this project's headline output is a predicted-versus-realised risk comparison, which such a leak would flatter exactly where the model is being judged.
  \item The regressors are the targets' own orthogonalised series, so the hierarchy compounds: by the time a level-3 factor is residualised, the level-0 factors it sees have already had everything above them removed.
  \item Before the first refit, with fewer than 252 usable observations or a missing regressor on the day, the orthogonalised value is NULL rather than a silent fallback to the raw value. A factor that quietly stops being orthogonal on some dates is worse than one with a documented gap.
\end{enumerate}

Two alternative modes exist and are not used. \code{expanding} fits on all history to date: causal, but it never forgets, so a loading that was high in one crisis stays high for a decade and the residual acquires a systematic negative loading on its own regressor. \code{full\_sample} fits once on everything: exactly orthogonal, and what commercial risk models do, but every historical factor value then contains information from its own future. Measured residual correlation of \code{eq\_em} against \code{eq\_global} over 2009--2026: rolling $-0.05$, expanding $-0.34$, full sample $0.00$.
